\chapter{المؤثرات المتراصة والمبرهنة الطيفية}
\label{ch:b3:spectral}

القُطرنة درّة تاج الجبر الخطي المنتهي البُعد: إذ لمصفوفة
متناظرة أساسٌ متعامد متجانس من المتجهات الذاتية. وفي البُعد
اللانهائي يخفق هذا من أجل المؤثرات المحدودة ذاتية المرافقة
عمومًا --- فالضرب في $x$ على $L^2(\intcc01)$ ليست له \emph{أي}
قيمة ذاتية البتة (\cref{exo:b3:spectral:6}) --- لكنه ينجو، في صورة تامة، من
أجل المؤثرات التي هي \emph{شبه منتهية البُعد}: أي المتراصة.
والمبرهنة الطيفية للمؤثرات المتراصة ذاتية المرافقة أكثرُ
مبرهنات التحليل الدالي التطبيقي استعمالًا على الإطلاق: فهي
تقطرن المعادلات التكاملية، وتقود \omterm{thm:b3:spectral:fredholm}{بديل فريدهولم}، و(في مسألة
نهاية الأسبوع) تحل الوتر المهتز، منتجةً أساسَ الجيوب في تحليل
فورييه من نظرية المؤثرات المحضة --- مع سقوط مجموع أولير $\zeta(2) = \frac{\pi^2}6$
من صيغة أثرٍ هديةً عند الوداع. وفي كل ما يلي، يكون $H$ \omterm{def:b3:hilbert:inner}{فضاء
هيلبرت} على $\C$ (أو على $\R$؛ وتتكيّف العبارات)، وتكون
المؤثرات محدودة.

\section{المؤثرات المتراصة}

\begin{definition}\label{def:b3:spectral:compact}
يكون $T \in \mathcal L(E, F)$ (حيث $E, F$ فضاءا باناخ) \emph{متراصًّا}\index{مؤثر
متراص} إذا كانت صورة الكرة الواحدية $T(B)$ متراصةً نسبيًّا في
$F$ --- أي، بالتكافؤ، إذا كان لكل متتالية محدودة $(x_n)$
متتاليةٌ جزئية تجعل $(Tx_{n_k})$ متقاربة. والمؤثرات المنتهية
الرتبة متراصة (لأن المجموعات المحدودة في البُعد المنتهي كذلك)؛
ومطابقة فضاء لانهائي البُعد ليست متراصة أبدًا (مبرهنة ريس،
السنة الجامعية 2).
\end{definition}

\begin{proposition}\label{prop:b3:spectral:ideal}
تشكّل المؤثرات المتراصة $\mathcal K(E, F)$ فضاءً جزئيًّا مغلقًا من $\mathcal L(E, F)$،
ومثاليًّا ذا جانبين: فإذا كان $S$ \omterm{def:b3:spectral:compact}{متراصًّا} $\Rightarrow$ كان
$AS$ و $SB$ متراصَّين من أجل $A, B$ محدودين. وأكثر من ذلك، في
\omterm{def:b3:hilbert:inner}{فضاء هيلبرت}، يكون كل \omterm{def:b3:spectral:compact}{مؤثر متراص} نهايةً بالمعيار لمؤثرات منتهية
الرتبة.
\end{proposition}

\begin{proof}
الفضاء الجزئي: واضح من تمييز المتتاليات. \omterm{def:b3:rings:ideal}{والمثالي}: إذ ترسل
التطبيقات المحدودة المتتالياتِ المتقاربة إلى متقاربة والمحدودة
إلى محدودة. والإغلاق: ليكن $T_n \to T$ مع $T_n$ متراصة، ولتكن
$(x_k)$ محدودة بالعدد $1$؛ يجعل استخراج قطري
$(T_nx_{k_j})_j$ متقاربة من أجل كل $n$؛ عندئذٍ تكون
$(Tx_{k_j})$ كوشية، لأن
\[
\norm{Tx_{k_j} - Tx_{k_l}} \leq 2\vertiii{T - T_n} +
\norm{T_nx_{k_j} - T_nx_{k_l}} ,
\]
باختيار $n$ أولًا ثم الأدلّة. وأما التقريب في فضاءات هيلبرت:
فليكن $T$ \omterm{def:b3:spectral:compact}{متراصًّا}، و $K = \overline{T(B)}$ متراصة؛ وبإعطاء $\varepsilon$،
نغطي $K$ بعدد منتهٍ من الكرات $B(y_i, \varepsilon)$ وليكن $P$ الإسقاطَ
المتعامد على $V = \operatorname{Vect}(y_1, \dots, y_m)$ (وهو مغلق: لأنه منتهي البُعد). عندئذٍ يكون
للمؤثر $PT$ رتبةٌ منتهية، ومن أجل $\norm x \leq 1$: باختيار $y_i$
تحقق $\norm{Tx - y_i} <
\varepsilon$،
\[
\norm{Tx - PTx} \leq \norm{Tx - y_i} + \norm{P(y_i - Tx)}
\leq 2\varepsilon
\]
($y_i = Py_i$؛ $\vertiii P \leq 1$): ومنه $\vertiii{T - PT} \leq
2\varepsilon$.
\end{proof}

\begin{example}\label{ex:b3:spectral:examples}
(a) المؤثرات القطرية على $\ell^2$: يكون $T(x_n) = (\lambda_nx_n)$ \omterm{def:b3:spectral:compact}{متراصًّا} إذا
وفقط إذا كان $\lambda_n \to 0$ (\cref{exo:b3:spectral:2}).
(b) مؤثرات النواة على $\mathcal C(\intcc01)$: متراصة بأسكولي (\cref{exo:b3:complete:7}).
(c) \emph{مؤثرات هيلبرت--شميدت}\index{مؤثر هيلبرت--شميدت}: من
أجل $k \in L^2(\intcc01^2)$،
\[
(T_kf)(x) = \int_0^1k(x, y)\,f(y)\,\dd y
\]
يعرّف مؤثرًا \omterm{def:b3:spectral:compact}{متراصًّا} على $L^2(\intcc01)$ مع
$\vertiii{T_k} \leq \norm k_{L^2}$
(\cref{exo:b3:spectral:4}: إذ يُبرز بترُ نشر $k$ في الأساس المؤثرَ $T_k$
نهايةً لمؤثرات منتهية الرتبة).
\end{example}

\section{المؤثرات ذاتية المرافقة}

\begin{definition}\label{def:b3:spectral:selfadjoint}
يكون $T \in \mathcal L(H)$ \emph{ذاتي المرافقة}\index{مؤثر ذاتي المرافقة} إذا
كان $T = T^*$ (\cref{exo:b3:hilbert:8})، أي $\langle Tx, y\rangle = \langle x, Ty\rangle$ من أجل كل $x, y$. عندئذٍ
يكون $\langle x, Tx\rangle \in \R$ من أجل كل $x$ (لأنه يساوي مرافقه).
\end{definition}

\begin{proposition}\label{prop:b3:spectral:sanorm}
من أجل $T$ \omterm{def:b3:spectral:selfadjoint}{ذاتي المرافقة}:\index{نصف القطر العددي}
\[
\vertiii T = \sup_{\norm x \leq 1}\ \abs{\langle x,
Tx\rangle} .
\]
والقيم الذاتية للمؤثر $T$ حقيقية، والمتجهات الذاتية الموافقة
لقيم ذاتية متمايزة متعامدة.
\end{proposition}

\begin{proof}
ليكن $M$ السوپريموم؛ فإن $M \leq \vertiii T$ بكوشي--شوارتز. وبالعكس، تعطي
المتطابقة من نمط الاستقطاب
\[
\langle x{+}y, T(x{+}y)\rangle - \langle x{-}y,
T(x{-}y)\rangle = 4\operatorname{Re}\langle y, Tx\rangle
\]
(بالنشر؛ إذ تتلاشى الحدود المتقاطعة $\langle y, Tx\rangle + \langle x,
Ty\rangle = 2\operatorname{Re}\langle y, Tx\rangle$ بذاتية المرافقة)،
مع قاعدة متوازي الأضلاع،
\[
4\operatorname{Re}\langle y, Tx\rangle \leq
M\bigl(\norm{x{+}y}^2 + \norm{x{-}y}^2\bigr) = 2M\bigl(\norm
x^2 + \norm y^2\bigr).
\]
ومن أجل $\norm x = 1$ مع $Tx \neq 0$، نأخذ $y = Tx/\norm{Tx}$: $4\norm{Tx} \leq 4M$. ومنه
$\vertiii T \leq M$. وأما القيم الذاتية: فإن $Tx = \lambda x$ مع $x \ne 0$
يعطي $\lambda\norm x^2 = \langle
x, Tx\rangle \in \R$. وأما التعامد: $\lambda\langle x,
y\rangle = \langle Tx, y\rangle = \langle x, Ty\rangle =
\mu\langle x, y\rangle$ مع $\lambda \neq \mu$ حقيقيتين.
\end{proof}

\section{المبرهنة الطيفية}

\begin{lemma}[وجود قيمة ذاتية قصوى]
\label{lem:b3:spectral:existence}
ليكن $T \neq 0$ \omterm{def:b3:spectral:compact}{متراصًّا} \omterm{def:b3:spectral:selfadjoint}{وذاتي المرافقة}. عندئذٍ يكون
$\vertiii T$ أو $-\vertiii T$ قيمةً ذاتية للمؤثر $T$.
\end{lemma}

\begin{proof}
حسب \cref{prop:b3:spectral:sanorm}، نختار متجهات واحدية $x_n$ تحقق $\langle x_n, Tx_n\rangle \to \mu$، حيث
$\abs\mu =
\vertiii T > 0$ (وننتقل إلى متتالية جزئية لتثبيت الإشارة). عندئذٍ
\[
\norm{Tx_n - \mu x_n}^2
= \norm{Tx_n}^2 - 2\mu\langle x_n, Tx_n\rangle + \mu^2
\leq \vertiii T^2 - 2\mu\langle x_n, Tx_n\rangle + \mu^2
\longrightarrow 2\mu^2 - 2\mu\cdot\mu = 0 .
\]
وبالتراص، تكون متتالية جزئية $Tx_{n_k} \to y$؛ عندئذٍ $\mu
x_{n_k} = Tx_{n_k} - (Tx_{n_k} - \mu x_{n_k}) \to y$،
ومنه $x_{n_k} \to x = y/\mu$، وهي متجهة واحدية، ويعطي الاتصال $Tx = \mu x$.
\end{proof}

\begin{theorem}[المبرهنة الطيفية للمؤثرات المتراصة ذاتية
المرافقة]\label{thm:b3:spectral:spectral}
ليكن $T$ مؤثرًا \omterm{def:b3:spectral:compact}{متراصًّا} \omterm{def:b3:spectral:selfadjoint}{ذاتي المرافقة} على \omterm{def:b3:hilbert:inner}{فضاء هيلبرت}
$H$.\index{المبرهنة الطيفية (المتراصة ذاتية المرافقة)}
\begin{enumerate}
  \item يقبل $H$ نظامًا متعامدًا متجانسًا $(e_n)_{n \in N}$ (حيث $N$
        منتهٍ أو قابل للعدّ) من المتجهات الذاتية للمؤثر $T$، بقيم
        ذاتية حقيقية غير معدومة $(\lambda_n)$، بحيث
        \[
        Tx = \sum_{n\in N}\lambda_n\,\langle e_n, x\rangle\,
        e_n \qquad (x \in H),
        \]
        و $H = \ker T \,\oplus^\perp\,
        \overline{\operatorname{Vect}}(e_n : n \in N)$.
  \item وإذا كان $N$ لانهائيًّا، فإن $\lambda_n \to 0$؛ ومن أجل كل
        $\delta > 0$ لا يحقق إلا عددٌ منتهٍ من $n$ الشرطَ
        $\abs{\lambda_n} \geq \delta$، ويكون كل فضاء ذاتي $\ker(T - \lambda)$، حيث
        $\lambda \neq 0$، منتهي البُعد.
  \item وبإكمال $(e_n)$ بأساس متعامد متجانس للنواة $\ker T$
        نحصل، حين يكون $H$ قابلًا للفصل، على أساس متعامد
        متجانس للفضاء $H$ مؤلَّف من متجهات ذاتية: أي إن $T$
        مقطرن.
\end{enumerate}
\end{theorem}

\begin{proof}
نبدأ بالجزء (2). لو كان لعدد لانهائي من المتجهات الذاتية المتعامدة
المتجانسة $x_k$ أن يحقق $\abs{\lambda_{(k)}} \geq \delta$: لكان $\norm{Tx_k - Tx_l}^2 =
\lambda_{(k)}^2 + \lambda_{(l)}^2 \geq 2\delta^2$ (بالتعامد
وفيثاغورس): فلا توجد أي متتالية جزئية متقاربة من $(Tx_k)$،
وهو يناقض تراص $T$ على المتتالية المحدودة $(x_k)$. وهذا يحدّ
بعدد منتهٍ، من أجل كل $\delta$، الترافقَ الكلي للقيم الذاتية
خارج $\intoo{-\delta}\delta$؛ وينتج من ذلك العدُّ و $\lambda_n \to 0$.

(1) ليكن $H_0$ الفضاء المولَّد المغلق بواسطة \emph{جميع} المتجهات
الذاتية ذات القيم الذاتية غير المعدومة، منظَّمةً (حسب (2)
وغرام--شميدت داخل كل فضاء ذاتي منتهي البُعد، مع التعامد بين
الفضاءات الذاتية من \cref{prop:b3:spectral:sanorm}) في نظام متعامد متجانس $(e_n)$
بقيم ذاتية $\lambda_n \neq 0$. ويرسل $T$ الفضاءَ $H_0$ في $H_0$، كما يرسل
$H_0^\perp$ في $H_0^\perp$: إذ من أجل $y
\perp H_0$ ومن أجل $e$ متجهة
ذاتية، $\langle e, Ty\rangle =
\langle Te, y\rangle = \lambda\langle e, y\rangle = 0$. والتضييق $T' = T\restriction_{H_0^\perp}$ متراص \omterm{def:b3:spectral:selfadjoint}{وذاتي المرافقة} على \omterm{def:b3:hilbert:inner}{فضاء
هيلبرت} $H_0^\perp$؛ فإذا كان $T' \neq
0$، أنتجت \cref{lem:b3:spectral:existence} متجهةً
ذاتية للمؤثر $T$ بقيمة ذاتية غير معدومة داخل $H_0^\perp$ ---
وهو مستحيل، إذ تسكن هذه المتجهات في $H_0$. ومنه $T' = 0$: أي
$H_0^\perp \subseteq \ker T$. وبالعكس $\ker T \perp$ كل $e_n$ ($\langle e_n, z\rangle = \frac1{\lambda_n}\langle
Te_n, z\rangle = \frac1{\lambda_n}\langle e_n, Tz\rangle =
0$): أي
$\ker T \subseteq H_0^\perp$، ومنه $\ker T =
H_0^\perp$ والتفكيك المتعامد. وأما النشر: فمن أجل
$x = z + \sum_nc_ne_n$ (حيث $z \in \ker T$ و $c_n = \langle
e_n, x\rangle$؛ \cref{thm:b3:hilbert:parseval}(1) على $H_0$)،
يعطي اتصال $T$ أن $Tx = \sum_nc_n\lambda_ne_n$.

(3) النواة $\ker T$، وهي فضاء جزئي مغلق من فضاء قابل للفصل،
قابلةٌ للفصل: فلها أساس متعامد متجانس (\cref{prop:b3:hilbert:gramschmidt})؛ والاتحاد
أساسٌ متعامد متجانس للفضاء $H$ بحكم التفكيك في (1).
\end{proof}

\begin{theorem}[بديل فريدهولم]
\label{thm:b3:spectral:fredholm}
ليكن $T$ \omterm{def:b3:spectral:compact}{متراصًّا} \omterm{def:b3:spectral:selfadjoint}{ذاتي المرافقة} وليكن $\lambda \in \R\setminus
\{0\}$.\index{بديل
فريدهولم}
\begin{enumerate}
  \item إذا لم يكن $\lambda$ قيمةً ذاتية، فإن $T - \lambda I$
        تقابلٌ ذو مقلوب محدود: أي إن للمعادلة $Tx - \lambda x = f$، من أجل
        كل $f$، حلًّا واحدًا بالضبط يتعلق اتصاليًّا بالدالة $f$.
  \item وإذا كان $\lambda$ قيمةً ذاتية، فإن $Tx - \lambda x = f$ \omterm{def:b3:groups:derived}{قابلة للحل}
        إذا وفقط إذا كان $f \perp \ker(T - \lambda I)$، ويكون الحل وحيدًا إلى غاية
        تلك النواة (المنتهية البُعد).
\end{enumerate}
\end{theorem}

\begin{proof}
نفكّك $x = z + \sum c_ne_n$ و $f = w + \sum d_ne_n$ على امتداد \cref{thm:b3:spectral:spectral} ($z, w \in \ker T$). فتُقرأ
المعادلة
\[
-\lambda z = w, \qquad (\lambda_n - \lambda)\,c_n = d_n\
(n \in N).
\]
(1) $\lambda \notin \{\lambda_n\}\cup\{0\}$: فحسب (2) من المبرهنة الطيفية، $\inf_n\abs{\lambda_n - \lambda} = \delta >
0$ (إذ لا تتراكم
القيم الذاتية إلا عند $0 \neq \lambda$). ونحل: $z = -w/\lambda$
و $c_n = d_n/(\lambda_n - \lambda)$، مع $\sum\abs{c_n}^2 \leq \delta^{-2}\sum\abs{d_n}^2$: أي حلٌّ وحيد يحقق $\norm x \leq C\norm f$.
(2) $\lambda = \lambda_{n}$ من أجل $n$ في مجموعة منتهية $F$: فتقتضي قابلية
$(\lambda_n - \lambda)c_n = d_n$ للحل من أجل $n \in F$ أن يكون $d_n = 0$، أي
$f \perp e_n$ (حيث $n \in F$)، أي $f
\perp \ker(T - \lambda I)$؛ وتكون المعاملات $c_n$
من أجل $n \in F$ حرةً عندئذٍ.
\end{proof}

\begin{example}\label{ex:b3:spectral:minxy}
على $L^2(\intcc01)$، لتكن $Tf(x) = \int_0^1\min(x, y)f(y)\dd y$: وهي \omterm{ex:b3:spectral:examples}{مؤثر هيلبرت--شميدت} ذو
نواة حقيقية متناظرة: فهو متراص \omterm{def:b3:spectral:selfadjoint}{وذاتي المرافقة}. وبحل
$Tf = \lambda f$: تبيّن العلاقة $\bigl(Tf\bigr)(x) = \int_0^xyf(y)\dd y +
x\int_x^1f(y)\dd y$ أن $u = Tf$ يحقق
$u'' = -f$ (باشتقاقين، وهو مشروع من أجل $f$ متصلة، و $Tf$ متصلة
من أجل $f \in L^2$: بالتقارب المهيمن)، مع $u(0)
= 0$ و $u'(1) = 0$.
ومنه تحل الدوال الذاتية $\lambda u'' =
-u$ مع $u(0) = 0$ و $u'(1) = 0$:
\[
u_n(x) = \sin\Bigl(\bigl(n + \tfrac12\bigr)\pi x\Bigr),
\qquad
\lambda_n = \frac{1}{\bigl(n + \frac12\bigr)^2\pi^2}
\quad (n \geq 0),
\]
وتؤكد المبرهنة الطيفية --- دون أي نظرية فورييه --- أن هذه
الجيوب تشكّل، بعد التعيير، أساسًا متعامدًا متجانسًا للفضاء
$L^2(\intcc01)$ (إذ نواة $T$ هي $0$: إذ إن $Tf = 0$ يفرض،
بالاشتقاقين، أن $f = 0$ في كل مكان تقريبًا). وتجري مسألة نهاية
الأسبوع دائرة الأفكار نفسها من أجل الوتر المهتز وتستخرج
$\zeta(2)$ من الأثر.
\end{example}

\begin{method}\label{met:b3:spectral:method}
بإعطاء معادلة تكاملية أو تفاضلية: (1) أعد صياغتها على الصورة
$(I - \lambda K)u = f$ أو $Ku = \lambda u$ حيث $K$ مؤثر تكاملي؛ (2) تحقق من
تراص $K$ (بنواة هيلبرت--شميدت، أو بأسكولي) ومن ذاتية مرافقته
إن أمكن (بنواة حقيقية متناظرة)؛ (3) قطرِنه بالمبرهنة الطيفية أو
استدعِ \omterm{thm:b3:spectral:fredholm}{بديل فريدهولم} للقابلية للحل؛ (4) واقرأ الوجود
والوحدانية والاستقرار وصيغ المتسلسلات للحلول في الأساس الذاتي.
والمؤثرات التفاضلية غير محدودة، لكن \emph{مقلوباتها} (أي مؤثرات
غرين) متراصة: فاقلب أولًا دائمًا.
\end{method}

\section{تمارين}

\begin{exercise}[$\star$]\label{exo:b3:spectral:1}
(a) برهن على أن مؤثرًا محدودًا مجالُه منتهي البُعد متراصٌّ.
(b) برهن على أن مطابقة فضاء معياري متراصة إذا وفقط إذا كان
البُعد منتهيًا (ريس، السنة الجامعية 2). واستنتج أن مؤثرًا
\omterm{def:b3:spectral:compact}{متراصًّا} على فضاء لانهائي البُعد لا يكون أبدًا قابلًا للقلب
بمقلوب محدود.
\end{exercise}

\begin{exercise}[$\star$]\label{exo:b3:spectral:2}
لتكن $T(x_1, x_2, \dots) = (\lambda_1x_1, \lambda_2x_2, \dots)$ على $\ell^2$، مع $(\lambda_n)$ محدودة.
(a) برهن على $\vertiii T = \sup\abs{\lambda_n}$.
(b) برهن على أن $T$ متراص إذا وفقط إذا كان $\lambda_n \to 0$.
\emph{(من أجل $\Leftarrow$، ابتر؛ ومن أجل $\Rightarrow$، اختبر
على $(e_n)$.)}
(c) متى يكون $T$ \omterm{def:b3:spectral:selfadjoint}{ذاتي المرافقة}؟ وتحقق من المبرهنة الطيفية
بالمعاينة في تلك الحالة.
\end{exercise}

\begin{exercise}[$\star\star$]\label{exo:b3:spectral:3}
أعطِ تفاصيل خاصية \omterm{def:b3:rings:ideal}{المثالي} (\cref{prop:b3:spectral:ideal}): أنه إذا كان $S$
\omterm{def:b3:spectral:compact}{متراصًّا} وكان $A, B$ محدودين، فإن $ASB$ متراص. واستنتج أنه إذا
كان $ST = TS = I$ من أجل $S$ محدود ما، وكان $\dim H = \infty$، فإن $T$
ليس \omterm{def:b3:spectral:compact}{متراصًّا} --- ووفّق بين ذلك وبين \cref{exo:b3:spectral:1}(b).
\end{exercise}

\begin{exercise}[$\star\star$]\label{exo:b3:spectral:4}
(هيلبرت--شميدت) لتكن $k \in L^2(\intcc01^2)$ ولتكن $(e_n)$ \omterm{def:b3:hilbert:onb}{أساسًا هيلبرتيًّا}
للفضاء $L^2(\intcc01)$.
(a) برهن على أن $\vertiii{T_k} \leq \norm k_{L^2}$ \emph{(بكوشي--شوارتز في المتغيّر $y$،
ثم بتونيلي)}.
(b) انشر $k(x,y) = \sum_{m,n}c_{mn}e_m(x)\overline{e_n(y)}$ في $L^2$ للمربّع (وبرّر أن الجداءات تشكّل
\omterm{def:b3:hilbert:onb}{أساسًا هيلبرتيًّا} هناك)، وبرهن على أن بتر المجموع يعطي مؤثرات
منتهية الرتبة تتقارب إلى $T_k$ بمعيار المؤثرات: أي إن $T_k$
متراص.
\end{exercise}

\begin{exercise}[$\star\star$]\label{exo:b3:spectral:5}
ليكن $T$ \omterm{def:b3:spectral:selfadjoint}{ذاتي المرافقة} مع $\langle x, Tx\rangle \geq 0$ من أجل كل $x$ (أي مؤثر
\emph{موجب}).
(a) برهن على أن القيم الذاتية $\geq 0$ وعلى أن $\vertiii T =
\sup_{\norm x\leq1}\langle x, Tx\rangle$.
(b) برهن على متراجحة كوشي--شوارتز المعمَّمة $\abs{\langle x, Ty\rangle}^2 \leq \langle x, Tx\rangle\langle
y, Ty\rangle$.
\end{exercise}

\begin{exercise}[$\star\star$]\label{exo:b3:spectral:6}
على $L^2(\intcc01)$، ليكن $(Mf)(x) = x\,f(x)$.
(a) برهن على أن $M$ محدود \omterm{def:b3:spectral:selfadjoint}{وذاتي المرافقة} مع $\vertiii M =
1$، لكن ليست
له \emph{أي} قيمة ذاتية.
(b) برهن على أن $M$ ليس \omterm{def:b3:spectral:compact}{متراصًّا} \emph{(أبرز متتالية محدودة
ليس لصورتها أي متتالية جزئية متقاربة، مثلًا الدوال المميّزة
المعيَّرة لفترات متقلّصة بجوار $1$ --- أو استدعِ المبرهنة
الطيفية)}.
(c) وأين ينكسر برهان \cref{lem:b3:spectral:existence} من أجل $M$؟
\end{exercise}

\begin{exercise}[$\star\star$]\label{exo:b3:spectral:7}
(فولتيرا) على $L^2(\intcc01)$، ليكن $Vf(x) = \int_0^xf(y)\dd
y$.
(a) برهن على أن $V$ متراص (فهو هيلبرت--شميدت بالنواة $\mathbf
1_{y < x}$)
لكنه ليس \omterm{def:b3:spectral:selfadjoint}{ذاتي المرافقة}؛ واحسب $V^*$.
(b) برهن على أنه ليست للمؤثر $V$ أي قيمة ذاتية غير معدومة.
\emph{(انطلاقًا من $Vf =
\lambda f$: تقبل $f$ ممثّلًا \omterm{def:b3:topology:continuity}{متصلًا}، ثم تكون
من الصنف $\mathcal C^1$، وتحل $\lambda f' = f$ مع
$f(0) = 0$.)}
(c) اخلص إلى أن التراص وحده لا يعطي أي متجهات ذاتية: فذاتية
المرافقة في \cref{thm:b3:spectral:spectral} جوهرية.
\end{exercise}

\begin{exercise}[$\star\star\star$]\label{exo:b3:spectral:8}
(كورانت--فيشر) ليكن $T$ \omterm{def:b3:spectral:compact}{متراصًّا} \omterm{def:b3:spectral:selfadjoint}{وذاتي المرافقة} و\emph{موجبًا}،
بقيم ذاتية $\mu_1 \geq \mu_2 \geq \dots
> 0$ (مكرَّرة حسب الترافق، بمتجهات ذاتية
$e_1, e_2,
\dots$). برهن على:
\[
\mu_{k} = \max_{\substack{V \subseteq H \\ \dim V = k}}\
\min_{\substack{x \in V\\ \norm x = 1}}\ \langle x, Tx\rangle
= \min_{\substack{W \subseteq H\\ \operatorname{codim}W =
k-1}}\ \max_{\substack{x\in W\\ \norm x = 1}}\ \langle x,
Tx\rangle .
\]
\emph{(اختبر $V = \operatorname{Vect}(e_1,\dots,e_k)$؛ ومن أجل الحدّ الأعلى قاطع أي $V$ مع
فضاءات من نمط $\operatorname{Vect}(e_k, e_{k+1}, \dots)$: فيفرض عدُّ الأبعاد تقاطعًا غير معدوم.)}
واستنتج أن القيم الذاتية تتعلق رتيبًا بالمؤثر $T$ ($T \leq S
\Rightarrow \mu_k(T) \leq \mu_k(S)$).
\end{exercise}

\begin{exercise}[$\star\star$]\label{exo:b3:spectral:9}
باستعمال \cref{thm:b3:spectral:fredholm} من أجل $Tf(x) =
\int_0^1\min(x,y)f(y)\dd y$ (\cref{ex:b3:spectral:minxy}): من أجل أي
$\lambda \in \R$ يكون للمعادلة التكاملية
\[
f(x) - \lambda\int_0^1\min(x,y)\,f(y)\,\dd y = g(x)
\]
حلٌّ وحيد $f \in L^2$ من أجل كل $g \in L^2$؟ وماذا يحدث عند
القيم الاستثنائية؟
\end{exercise}

\begin{exercise}[$\star\star\star$]\label{exo:b3:spectral:10}
ليكن $S$ الإزاحةَ على $\ell^2$ (\cref{exo:b3:banach:1}).
(a) برهن على أنه ليست للمؤثر $S$ أي قيمة ذاتية، بينما كل
$\lambda$ يحقق $\abs\lambda < 1$ قيمةٌ ذاتية للمؤثر $S^*$ (وجد المتجهات
الذاتية صراحةً، وهي متتاليات هندسية).
(b) وليس $S$ ولا $S^*$ \omterm{def:b3:spectral:compact}{متراصًّا}: تحقق من ذلك باختبار على
$(e_n)$ على غرار \cref{exo:b3:spectral:2}.
(c) علّق: أن مشهد القيم الذاتية، من أجل المؤثرات غير ذاتية
المرافقة وغير المتراصة، قد يكون أي شيء من الخالي إلى قرص كامل
--- والمفهوم الذي ينجو هو \emph{الطيف}، وهو مدروس في مقرر
لاحق.
\end{exercise}

\begin{exercise}[$\star\star$]\label{exo:b3:spectral:11}
(الجذور التربيعية) ليكن $T$ \omterm{def:b3:spectral:compact}{متراصًّا} \omterm{def:b3:spectral:selfadjoint}{وذاتي المرافقة} وموجبًا
($\langle x, Tx\rangle \geq 0$) على \omterm{def:b3:hilbert:inner}{فضاء هيلبرت} $H$، بتفكيك طيفي $Tx = \sum_n\mu_n\langle e_n,
x\rangle e_n$ (حيث
$\mu_n > 0$).
(a) عرّف $Sx = \sum_n\sqrt{\mu_n}\,\langle e_n, x\rangle
e_n$؛ وبرهن على أن $S$ متراص \omterm{def:b3:spectral:selfadjoint}{وذاتي المرافقة} وموجب
مع $S^2
= T$.
(b) برهن على الوحدانية: أن كل \omterm{def:b3:spectral:compact}{مؤثر متراص} موجب \omterm{def:b3:spectral:selfadjoint}{ذاتي المرافقة} $R$
يحقق $R^2 = T$ يحفظ الفضاءات الذاتية للمؤثر $T$
\emph{(إذ $RT = R^3 = TR$: ومنه يتبادل $R$ مع $T$، ومنه
$R(\ker(T - \mu)) \subseteq \ker(T - \mu)$)}، وعلى $\ker(T - \mu)$ يكون $R$ مؤثرًا موجبًا مربّعه
$\mu\,\mathrm{id}$ على فضاء منتهي البُعد: فقطرِنه هناك واخلص إلى $R =
\sqrt\mu\,\mathrm{id}$
على كل فضاء ذاتي، ومنه $R = S$.
(c) احسب $\sqrt G$ من أجل مؤثر الوتر $G$ في \cref{pb:b3:spectral:1}: أي نواة
لها القيم الذاتية $\frac1{n\pi}$ في أساس الجيوب؟ (عبّر عن
$\sqrt G$ بوصفه نهاية نوى بالمعنى $L^2$؛ ولا تُطلب أي صيغة
مغلقة.)
\end{exercise}

\begin{exercise}[$\star\star\star$]\label{exo:b3:spectral:12}
(تفكيك القيم المفردة) ليكن $T \in \mathcal L(H)$ \emph{\omterm{def:b3:spectral:compact}{متراصًّا}}، وليس بالضرورة
\omterm{def:b3:spectral:selfadjoint}{ذاتي المرافقة}.
(a) برهن على أن $T^*T$ متراص \omterm{def:b3:spectral:selfadjoint}{وذاتي المرافقة} وموجب؛ ولتكن
$(e_n)$ عائلةً متعامدة متجانسة من المتجهات الذاتية تحقق
$T^*Te_n = s_n^2e_n$ مع $s_n > 0$ (أي \emph{القيم المفردة})، مكمَّلةً
بالفضاء $\ker(T^*T) = \ker T$ (وبرهن على هذا التساوي).
(b) ضع $f_n = \frac{Te_n}{s_n}$؛ وبرهن على أن $(f_n)$ متعامدة متجانسة، وأثبت
\emph{تفكيك القيم المفردة}:
\[
Tx = \sum_n s_n\,\langle e_n, x\rangle\,f_n
\qquad (x \in H),
\]
مع التقارب في $H$.
(c) استنتج: أن $\vertiii T = \max_ns_n$؛ وأن $T$ نهايةٌ بالمعيار لمؤثرات منتهية
الرتبة (فيُعاد البرهان على عكس \cref{prop:b3:spectral:ideal} من أجل فضاءات
هيلبرت)؛ ومن أجل مؤثر فولتيرا $V$ في \cref{exo:b3:spectral:7}، الذي ليست له
أي قيمة ذاتية، اشرح لماذا يوجد تفكيك القيم المفردة رغم ذلك وما
مكوّناته (وعيّن $V^*V$ بوصفه مؤثر نواة من نمط الوتر --- وحساب
قيمه الذاتية صراحةً من اختصاص \cref{exo:b3:spectral:9}).
\end{exercise}

\section{مسألة: الوتر المهتز و
\texorpdfstring{$\zeta(2)$}{زيتا(2)}}

\begin{problem}[{مسألة نهاية الأسبوع --- مؤثر غرين، وأساس
الجيوب، وصيغة أثر}]\label{pb:b3:spectral:1}
نحل مسألة القيم الذاتية للوتر المهتز ذي الطرفين المثبَّتين ---
$-u'' = \nu u$ مع $u(0) = u(1) = 0$ --- بنظرية المؤثرات، فنحصل على أساس
الجيوب المتعامد المتجانس دون أي حساب فورييهي، ونقيّم $\zeta(2)$
بمقارنة عبارتين عن \emph{أثر} مؤثر غرين. ونعرّف، على
$L^2(\intcc01)$،
\[
(Gf)(x) = \int_0^1 g(x,y)\,f(y)\,\dd y,
\qquad
g(x, y) = \min(x,y)\,\bigl(1 - \max(x,y)\bigr).
\]

\medskip
\noindent\textbf{الجزء الأول --- مؤثر غرين.}
\begin{enumerate}
  \item برهن على أن $g$ متصلة ومتناظرة مع $0 \leq g
        \leq \frac14$، وعلى أن $G$
        متراص \omterm{def:b3:spectral:selfadjoint}{وذاتي المرافقة} (\cref{ex:b3:spectral:examples}(c)).
  \item من أجل $f$ متصلة، برهن على أن $u = Gf$ من الصنف $\mathcal
        C^2$
        مع
        \[
        -u'' = f, \qquad u(0) = u(1) = 0
        \]
        \emph{(اكتب $u(x) = (1-x)\int_0^xyf(y)\dd y +
        x\int_x^1(1-y)f(y)\dd y$ واشتقّ مرتين)}. وبالعكس، إذا كان
        $u \in \mathcal C^2$ مع $u(0) = u(1)
        = 0$، فإن $G(-u'') = u$: أي إن $G$ يقلب
        مؤثر الوتر.
  \item برهن على أن $\ker G = \{0\}$ \emph{(إذا كان $Gf = 0$
        مع $f \in
        L^2$: اختبر مقابل $\varphi$ متصلة، وانقل $G$
        بالتناظر أو بفوبيني إلى $\varphi$، واستعمل المبرهنة
        المساعدة الأساسية \cref{cor:b3:lp:fundlemma} --- أو نظِّم)}، وعلى أن
        $G$ مؤثر موجب: $\langle f, Gf\rangle \geq 0$. \emph{(من أجل $f$ متصلة:
        $\langle f, Gf\rangle = \int_0^1 (u')^2$ مع $u
        = Gf$، بالتجزئة؛ ثم اخلص بالكثافة.)}
\end{enumerate}

\medskip
\noindent\textbf{الجزء الثاني --- القُطرنة: أساس الجيوب.}
\begin{enumerate}[resume]
  \item برهن على أن الدوال الذاتية للمؤثر $G$ بقيمة ذاتية
        $\lambda \ne 0$ هي، إلى غاية سلّميات، حلول $-\lambda u'' = u$ مع
        $u(0) = u(1) = 0$ \emph{(إذ تقبل الدالة الذاتية ممثّلًا \omterm{def:b3:topology:continuity}{متصلًا}
        --- لأن $Gf$ متصلة من أجل $f \in L^2$، ولماذا؟ ---
        ومنه تكون من الصنف $\mathcal C^2$ بترقية السؤال 2)}.
  \item حل مسألة الشروط الحدّية: فالقيم الذاتية للمؤثر $G$ هي
        $\lambda_n = \frac1{n^2\pi^2}$ (حيث $n \geq 1$)، بدوال ذاتية معيَّرة $e_n(x) =
        \sqrt2\,\sin(n\pi x)$؛
        وتحقق من التعامد والتجانس بالمكاملة المباشرة بوصفه
        فحصًا للسلامة.
  \item اخلص من \cref{thm:b3:spectral:spectral} ومن السؤال 3 إلى أن $\bigl(\sqrt2\sin(n\pi x)\bigr)_{n\geq1}$
        \emph{أساسٌ} متعامد متجانس للفضاء $L^2(\intcc01)$ ---
        دون حاجة إلى ستون--فايرشتراس ولا إلى متسلسلات فورييه.
        وانشر $f(x) = x(1-x)$ في هذا الأساس واكتب \omterm{thm:b3:hilbert:parseval}{بارسيفال} من
        أجلها.
\end{enumerate}

\medskip
\noindent\textbf{الجزء الثالث --- صيغة الأثر و $\zeta(2)$.}
\begin{enumerate}[resume]
  \item برهن على المتطابقتين
        \[
        \langle e_n, Ge_n\rangle = \lambda_n
        \quad\text{و}\quad
        \sum_{n\geq1}\lambda_n = \int_0^1 g(x,x)\,\dd x .
        \]
        وأما الثانية (أي \emph{صيغة الأثر}): فانشر
        $g(x, \cdot)$، من أجل $x$ مثبَّت، في الأساس $(e_n)$
        --- وبرهن على أن المعاملات هي $\lambda_ne_n(x)$، بحيث يكون
        $g(x, \cdot) = \sum_n\lambda_ne_n(x)\,e_n$ في $L^2$. وهنا تكون الجيوب صريحة: فتحقق
        مباشرةً من أن $\sum_n\lambda_ne_n(x)e_n(y)$ يتقارب \emph{بانتظام} على المربّع
        (بالمقارنة مع $\sum
        \frac2{n^2\pi^2}$)، ومنه يكون مجموعه \omterm{def:b3:topology:continuity}{متصلًا}،
        وبكونه ذا النشرات $L^2$ نفسها في $y$ من أجل كل $x$،
        يساوي $g(x,y)$ في كل مكان. ثم ضع $y = x$ وكامل حدًّا
        حدًّا.
  \item احسب $\int_0^1g(x,x)\dd x = \int_0^1x(1-x)\dd x =
        \frac16$، واخلص إلى
        \[
        \sum_{n\geq1}\frac{1}{n^2\pi^2} = \frac16,
        \qquad\text{أي}\qquad
        \boxed{\ \zeta(2) = \frac{\pi^2}{6}\ } :
        \]
        أي مجموع أولير من أثر مؤثر.
  \item استنبط $\zeta(2)$ بطريقة ثالثة: طبّق \omterm{thm:b3:hilbert:parseval}{بارسيفال} في أساس
        الجيوب على الدالة الثابتة $\mathbf 1$، واحسب $\sum_{n \text{ فردي}}\frac1{n^2}$،
        واخلص. ثم قارن الآليتين: بأي معنى تكون حجة الأثر في
        السؤالين 7 و8 «\omterm{thm:b3:hilbert:parseval}{بارسيفال} مطبَّقة على النواة كلها
        دفعةً واحدة»؟
\end{enumerate}

\medskip
\noindent\textbf{الجزء الرابع --- الوتر يهتز.}
\begin{enumerate}[resume]
  \item (فصل المتغيّرات، مركَّبًا) من أجل $f \in
        L^2$، نعرّف
        \[
        u(t, x) = \sum_{n\geq1}\;c_n\,
        \cos(n\pi t)\,\sqrt2\sin(n\pi x),
        \qquad c_n = \langle e_n, f\rangle .
        \]
        برهن على أن المتسلسلة تتقارب في $L^2(\intcc01)$ من أجل
        كل $t$، وعلى أن $t\mapsto u(t, \cdot)$ متصلة بقيم في $L^2$، وعلى أنه
        من أجل $f$ في الفضاء المولَّد بعدد منتهٍ من $e_n$ تحل
        معادلة الأمواج $\partial_t^2u =
        \partial_x^2u$ مع $u(0) = f$ و $\partial_tu(0) = 0$،
        بطرفين مثبَّتين. والقيم الذاتية $n^2\pi^2$ هي مربّعات
        الترددات: أي توافقيات الوتر --- اشرح التفسير الموسيقي
        للمبرهنة \cref{thm:b3:spectral:spectral} في فقرة واحدة.
\end{enumerate}

\medskip
\noindent\textbf{الجزء الخامس --- عوائد تغيّرية: طريقة القوى،
واستقرار فايل، وحدٌّ صارم على $\pi$.} ليكن $A$ مؤثرًا \omterm{def:b3:spectral:compact}{متراصًّا}
\omterm{def:b3:spectral:selfadjoint}{ذاتي المرافقة} \emph{موجبًا} بقيم ذاتية $\mu_1 \geq \mu_2 \geq \cdots > 0$ ومتجهات ذاتية
متعامدة متجانسة $(u_n)$؛ و $R(x) = \frac{\langle x,
Ax\rangle}{\norm x^2}$. وصيغ الأصغري--الأعظمي هي
\cref{exo:b3:spectral:8}؛ وهنا ننفقها.
\begin{enumerate}[resume]
  \item (طريقة القوى) من أجل $x \neq 0$ نكتب $m_p =
        \sum_n\mu_n^p\abs{\langle u_n, x\rangle}^2$. برهن على
        $m_pm_{p+2} \geq m_{p+1}^2$ (بكوشي--شوارتز)، واستنتج السلسلة
        \[
        R(x) \;\leq\; \frac{\langle Ax, Ax\rangle}{\langle x,
        Ax\rangle} \;\leq\; R(Ax) \;\leq\; \mu_1,
        \]
        وبرهن على أنه إذا كان $\langle u_1, x\rangle \neq 0$، فإن $R(A^kx) \to \mu_1$: أي إن تكرار
        المؤثر على أي متجهة عامة يحسب القيمة الذاتية العليا
        --- أي طريقة القوى في التحليل العددي، مصدَّقةً.
  \item (استقرار فايل) من أجل $A, B$ متراصَّين ذاتيَي المرافقة
        موجبَين، استنتج من \cref{exo:b3:spectral:8} أن
        \[
        \abs{\mu_n(A) - \mu_n(B)} \;\leq\; \vertiii{A - B}
        \qquad\text{من أجل كل} n :
        \]
        أي إن الطيف كله ليبشيتزي بالثابت $1$ بالنسبة إلى معيار
        المؤثرات --- فيمكن حساب القيم الذاتية لجمل متناظرة
        كبيرة من تقريبات بخطأ مضمون.
  \item طبّق حدّ رايلي على $G$ بدالة الاختبار $u(x) = x(1-x)$:
        حُلّ $-w'' = u$ مع $w(0) =
        w(1) = 0$ لتحصل على $Gu = w = \frac{x(1-x)(1 + x -
        x^2)}{12}$، واحسب
        \[
        \norm u_2^2 = \frac1{30}, \qquad
        \langle u, Gu\rangle = \frac1{12}\Bigl(\frac1{30} +
        \frac1{140}\Bigr) = \frac{17}{5040},
        \qquad R(u) = \frac{17}{168},
        \]
        واخلص إلى الحدّ \emph{الصارم} $\frac1{\pi^2}
        = \lambda_1 \geq \frac{17}{168}$، أي $\pi \leq
        \sqrt{168/17} < 3.1437$.
  \item خطوةٌ واحدة من سلسلة السؤال 11، على دالة الاختبار
        نفسها: باستعمال $\int_0^1(x - x^2)^4\dd x =
        \frac1{630}$، احسب
        \[
        \norm{Gu}_2^2 = \frac1{144}\Bigl(\frac1{30} +
        \frac2{140} + \frac1{630}\Bigr) = \frac{31}{90720},
        \qquad
        \frac{\langle Gu, Gu\rangle}{\langle u, Gu\rangle} =
        \frac{31}{306},
        \]
        واخلص إلى $\pi \leq \sqrt{306/31} < 3.1419$: أي تكاملان، وأربعة أرقام صحيحة.
        (وكل تكرار إضافي يربّع الدقة تقريبًا: إذ تقود فجوة
        المتجهات الذاتية $\lambda_1/\lambda_2 = 4$ التقاربَ الهندسي.)
\end{enumerate}

\medskip
\noindent\textbf{الجزء السادس --- أثر $G^2$ و $\zeta(4)$.}
\begin{enumerate}[resume]
  \item برهن على أن $\sum_n\lambda_n^2 =
        \iint_{\intcc01^2}g(x,y)^2\,\dd x\,\dd y$ \emph{(انشر $g$ في الأساس الجدائي
        $(e_m(x)e_n(y))_{m,n}$ للفضاء $L^2(\intcc01^2)$ --- وهو \omterm{def:b3:hilbert:onb}{أساس هيلبرتي}، قارن
        \cref{exo:b3:spectral:5} --- وطبّق \omterm{thm:b3:hilbert:parseval}{بارسيفال} على المربّع؛ ويعيّن
        السؤال 7 المعاملات)}.
  \item احسب التكامل المزدوج:
        \[
        \iint g^2 = 2\int_0^1(1-x)^2\Bigl(\int_0^x
        y^2\,\dd y\Bigr)\dd x
        = \frac23\int_0^1x^3(1-x)^2\,\dd x = \frac1{90} .
        \]
  \item اخلص إلى $\zeta(4) = \dfrac{\pi^4}{90}$؛ واشرح، دون حساب، كيف تنتج آثار القوى
        العليا $G^k$ المقاديرَ $\zeta(2k) \in \pi^{2k}\,\Q$ من أجل كل $k \geq 1$،
        ولماذا تكون القيم الفردية $\zeta(3),
        \zeta(5), \dots$ بعيدةً بنيويًّا عن
        متناول هذه الآلة.
  \item ($\pi$ من الأسفل) استنتج من $\lambda_1^2 \leq
        \sum_n\lambda_n^2 = \frac1{90}$ أن $\pi \geq
        90^{1/4} > 3.080$،
        واجمعه مع السؤال 14 لتحصل على الحكم ذي الطرفين
        \[
        3.080 \;<\; \pi \;<\; 3.1419,
        \]
        محصَّلًا عليه كليًّا من حساب الوتر المهتز. وأي الطرفين
        يتقارب أسرع إذا استُعملت آثار عليا $(\operatorname{tr}G^{2k})^{-1/4k}$، ولماذا؟
\end{enumerate}

\medskip
\noindent\textbf{الجزء السابع --- الإقحام والرنين.} نثبّت
$\nu \in \R$ وننظر في الوتر المُقحَم $-u'' - \nu u =
f$ مع $u(0) = u(1) = 0$، حيث
$f \in L^2$ و $c_n = \langle
e_n, f\rangle$.
\begin{enumerate}[resume]
  \item لنفترض أن $\nu \notin \{n^2\pi^2 : n \geq 1\}$. برهن على أن
        \[
        u = \sum_{n\geq1}\frac{c_n}{n^2\pi^2 - \nu}\,e_n
        \]
        تتقارب في $L^2$ وبانتظام على $\intcc01$
        \emph{(بكوشي--شوارتز بين $(c_n)$ وذيول $\sum n^{-4}$،
        مع $\norm{e_n}_\infty =
        \sqrt2$)}، وعلى أنها تحقق $u = Gf + \nu Gu$ --- وهي صورة \omterm{thm:b3:spectral:fredholm}{بديل
        فريدهولم} (\cref{thm:b3:spectral:fredholm}) بالإحداثيات الصريحة، مع
        الوحدانية.
  \item لنفترض أن $\nu = m^2\pi^2$. برهن على أن للمعادلة
        $u = Gf + \nu
        Gu$ حلًّا إذا وفقط إذا كان $c_m = 0$، وهو وحيد
        إلى غاية إضافة مضاعفات للدالة $e_m$. والقراءة
        الفيزيائية: دفع أرجوحة عند ترددها الخاص بالضبط.
  \item من أجل $\nu < \pi^2$، برهن على أن مؤثر الحل $R_\nu\colon f \mapsto u$
        محدود على $L^2$ بمعيار $\frac1{\pi^2 - \nu}$، ومتراص، \omterm{def:b3:spectral:selfadjoint}{وذاتي
        المرافقة}، وموجب: فيبدأ التحليل الطيفي كله من جديد،
        مزاحًا بمقدار $\nu$.
  \item (تركيب) اجمع معجم هذه المسألة: القيمة الذاتية $\leftrightarrow$
        مربّع التردد (التوافقيات)؛ والأثر $\leftrightarrow$ $\zeta(2)$؛
        ومعيار هيلبرت--شميدت $\leftrightarrow$ $\zeta(4)$؛ \omterm{thm:b3:spectral:fredholm}{وبديل فريدهولم}
        $\leftrightarrow$ الرنين؛ والأصغري--الأعظمي $\leftrightarrow$ الحدود
        التغيّرية ($\pi <
        3.1437$ من كثير حدود واحد). مؤثر تكاملي
        واحد، وخمسة فصول من التحليل مصروفة.
\end{enumerate}

\medskip
\noindent\textbf{الجزء الثامن --- ثلاثة أصداء أخيرة.}
\begin{enumerate}[resume]
  \item ($\zeta(6)$، مجانًا) أعطت متطابقة \omterm{thm:b3:hilbert:parseval}{بارسيفال} في السؤال 6
        المقدارَ $\sum_{n\text{ فردي}}n^{-6} =
        \frac{\pi^6}{960}$. اشطر $\zeta(6)$ إلى $n$ فردية وزوجية
        واخلص إلى
        \[
        \zeta(6) = \frac{\pi^6}{945},
        \]
        دون أي تكامل جديد: فكانت آلة السؤال 17 (أي آثار $G^3$)
        ستعطي القيمة نفسها بكلفة نواة مكرَّرة --- \omterm{thm:b3:hilbert:parseval}{وبارسيفال} على
        دالة واحدة حسنة الاختيار هي الطريق الأرخص هنا.
  \item (الحالة الأساسية موجبة) ليكن $A$ مؤثرًا \omterm{def:b3:spectral:compact}{متراصًّا} \omterm{def:b3:spectral:selfadjoint}{ذاتي
        المرافقة} موجبًا على $L^2(\intcc01)$ معطًى بنواة متصلة
        متناظرة $k > 0$ على $\intoo01^2$، بأكبر قيمة ذاتية
        $\mu_1$. برهن على: (a) أن كل معظِّم لنسبة رايلي دالةٌ
        ذاتية للقيمة $\mu_1$؛ (b) وأنه إذا كانت $u$ كذلك، فإن
        $\langle\abs u, A\abs u\rangle \geq \langle u,
        Au\rangle$، مع متراجحة \emph{أكيدة} إذا كانت $u$ تأخذ
        الإشارتين على مجموعات ذات \omterm{def:b3:measure:measure}{قياس} موجب --- ومنه تكون
        للدالة $u$ إشارة ثابتة في كل مكان تقريبًا، ولا ينعدم
        $u = \mu_1^{-1}Au$ أبدًا على $\intoo01$؛ (c) وأن $\mu_1$ قيمة
        ذاتية \emph{بسيطة}. وتحقق من كل ادعاء على $G$: $e_1 = \sqrt2\sin(\pi x) > 0$،
        وكل $e_n$ حيث $n
        \geq 2$، بكونها متعامدة مع $e_1$، يجب أن
        تغيّر إشارتها (وهي تفعل: لها $n - 1$ جذرًا داخليًّا).
  \item (المسافة إلى الطيف، وثمن الرنين) من أجل $\nu \notin \{n^2\pi^2\}$، برهن
        على أن مؤثر الحل $R_\nu$ في السؤال 19 محدود \omterm{def:b3:spectral:selfadjoint}{وذاتي
        المرافقة} ومتراص، مع
        \[
        \vertiii{R_\nu} =
        \frac1{\min_{n\geq1}\,\abs{n^2\pi^2 - \nu}}
        = \frac1{\operatorname{dist}\bigl(\nu,
        \{n^2\pi^2\}\bigr)},
        \]
        والمعيار مبلوغ على النمط الأقرب. ثم كمِّم أرجوحة السؤال
        20: فالإقحام بالدالة $f = e_1$ عند $\nu = (1 - \varepsilon)\pi^2$ ينتج $u =
        \frac{1}{\varepsilon\pi^2}\,e_1$،
        أي تضخيمٌ بالعامل $\frac1\varepsilon$ للاستجابة الساكنة $Ge_1 =
        \frac1{\pi^2}e_1$ ---
        فعند واحد في المئة تحت التردد الأساسي ($\varepsilon = 10^{-2}$)، يجيب
        الوتر أعلى بمئة مرة.
\end{enumerate}
\end{problem}
